% GENERATED FILE. Edit canonical lessons, then run npm run generate:books.
% canonical-source-hash: fnv1a64:8317327f3f9bea62
% canonical-lessons: MW-C53-khelno, MW-C53-bechno, MW-C53-deno, MW-C53-leno, MW-C53-rakhno

\chapter{To play, To sell, To give, To take, To keep}
\label{ch:mw-to-play-to-sell-to-give-to-take-to-keep}
\hlchaptermodality{53}

\begin{chapteropening}
\textbf{By the end of this chapter:} \emph{I can say to play, to sell, to give, to take and to keep, to put.}

Complete the last lesson of chapter 53: I can say to play, to sell, to give, to take and to keep, to put.
\end{chapteropening}

\section[\texorpdfstring{\mw{खेलणो} (khelṇo)}{खेलणो (khelṇo)}]{\mw{खेलणो} (khelṇo) — to play}

\label{lesson:MW-C53-khelno}



\begin{quote}
Before the new one: say the Marwadi for to walk, then the Marwadi for to run.
\end{quote}

\subsection*{You'll want to know: \mw{खेलणो}}
\textbf{\mw{खेलणो}} — \emph{khelṇo} — \textquotedblleft{}to play\textquotedblright{}.

\begin{cousinweb}
\textbf{\mw{खेलणो}} (\emph{khelṇo}) is a tadbhava, a word inherited and reshaped by sound change. It comes from \emph{*khel}, \textquotedblleft{}play\textquotedblright{}, a form CDIAL reconstructs from its descendants. CDIAL lists Old Marwari \emph{khelaï} under it (CDIAL 3918). CDIAL heads a group of reconstructed words for \textquotedblleft{}play\textquotedblright{} \emph{*khēḍ}, and says that if Sanskrit \emph{krīḍati}, \textquotedblleft{}plays\textquotedblright{}, is of non-Aryan origin, these may share its source.
\end{cousinweb}

\begin{grammarlens}[title={What you've built}]
One more word for this chapter.
\end{grammarlens}

\subsection*{Your turn}
\begin{itemize}
\raggedright
  \item \emph{Say it:} \emph{khelṇo}
  \item \emph{Say it:} \emph{khelṇo}, once more
  \item \emph{Say it:} say \emph{khelṇo}
  \item \emph{Recall:} say the Marwadi for to walk, then the Marwadi for to run, then say \emph{khelṇo} again
\end{itemize}

\begin{tcolorbox}[breakable,colback=teal!4,colframe=teal!35!black,title={Before you move on}]
What does \mw{खेलणो} mean? (\textquotedblleft{}To play\textquotedblright{}.) Say it once more.
\end{tcolorbox}
\section[\texorpdfstring{\mw{बेचणो} (bechṇo)}{बेचणो (bechṇo)}]{\mw{बेचणो} (bechṇo) — to sell}

\label{lesson:MW-C53-bechno}



\begin{quote}
Before the new one: say the Marwadi for to run, then the Marwadi for to play.
\end{quote}

\subsection*{You'll want to know: \mw{बेचणो}}
\textbf{\mw{बेचणो}} — \emph{bechṇo} — \textquotedblleft{}to sell\textquotedblright{}.

\begin{cousinweb}
\textbf{\mw{बेचणो}} (\emph{bechṇo}) is a tadbhava, a word inherited and reshaped by sound change. It comes from \emph{*vetyayati}, \textquotedblleft{}makes a price\textquotedblright{}, a form CDIAL reconstructs from its descendants. CDIAL gives no Marwari form in this entry (CDIAL 12100); the match is through Hindi \emph{becnā} and Gujarati \emph{vecvũ}. CDIAL builds it on Middle Indo-Aryan \emph{vēcā}, \textquotedblleft{}hire, wages\textquotedblright{}.
\end{cousinweb}

\begin{grammarlens}[title={What you've built}]
One more word for this chapter.
\end{grammarlens}

\subsection*{Your turn}
\begin{itemize}
\raggedright
  \item \emph{Say it:} \emph{bechṇo}
  \item \emph{Say it:} \emph{bechṇo}, once more
  \item \emph{Say it:} say \emph{bechṇo}
  \item \emph{Recall:} say the Marwadi for to run, then the Marwadi for to play, then say \emph{bechṇo} again
\end{itemize}

\begin{tcolorbox}[breakable,colback=teal!4,colframe=teal!35!black,title={Before you move on}]
What does \mw{बेचणो} mean? (\textquotedblleft{}To sell\textquotedblright{}.) Say it once more.
\end{tcolorbox}
\section[\texorpdfstring{\mw{देणो} (deṇo)}{देणो (deṇo)}]{\mw{देणो} (deṇo) — to give}

\label{lesson:MW-C53-deno}



\begin{quote}
Before the new one: say the Marwadi for to play, then the Marwadi for to sell.
\end{quote}

\subsection*{You'll want to know: \mw{देणो}}
\textbf{\mw{देणो}} — \emph{deṇo} — \textquotedblleft{}to give\textquotedblright{}.

\begin{cousinweb}
\textbf{\mw{देणो}} (\emph{deṇo}) is a tadbhava, a word inherited from Sanskrit and reshaped by sound change. It comes from \emph{dadāti}, \textquotedblleft{}gives\textquotedblright{}. CDIAL lists Marwari \emph{deṇo} under it (CDIAL 6141).
\end{cousinweb}

\begin{grammarlens}[title={What you've built}]
One more word for this chapter.
\end{grammarlens}

\subsection*{Your turn}
\begin{itemize}
\raggedright
  \item \emph{Say it:} \emph{deṇo}
  \item \emph{Say it:} \emph{deṇo}, once more
  \item \emph{Say it:} say \emph{deṇo}
  \item \emph{Recall:} say the Marwadi for to play, then the Marwadi for to sell, then say \emph{deṇo} again
\end{itemize}

\begin{tcolorbox}[breakable,colback=teal!4,colframe=teal!35!black,title={Before you move on}]
What does \mw{देणो} mean? (\textquotedblleft{}To give\textquotedblright{}.) Say it once more.
\end{tcolorbox}
\section[\texorpdfstring{\mw{लेणो} (leṇo)}{लेणो (leṇo)}]{\mw{लेणो} (leṇo) — to take}

\label{lesson:MW-C53-leno}



\begin{quote}
Before the new one: say the Marwadi for to sell, then the Marwadi for to give.
\end{quote}

\subsection*{You'll want to know: \mw{लेणो}}
\textbf{\mw{लेणो}} — \emph{leṇo} — \textquotedblleft{}to take\textquotedblright{}.

\begin{cousinweb}
\textbf{\mw{लेणो}} (\emph{leṇo}) is a tadbhava, a word inherited from Sanskrit and reshaped by sound change. It comes from \emph{labhate}, \textquotedblleft{}catches, takes\textquotedblright{}. CDIAL lists Old Marwari \emph{lei} under it (CDIAL 10948).
\end{cousinweb}

\begin{grammarlens}[title={What you've built}]
One more word for this chapter.
\end{grammarlens}

\subsection*{Your turn}
\begin{itemize}
\raggedright
  \item \emph{Say it:} \emph{leṇo}
  \item \emph{Say it:} \emph{leṇo}, once more
  \item \emph{Say it:} say \emph{leṇo}
  \item \emph{Recall:} say the Marwadi for to sell, then the Marwadi for to give, then say \emph{leṇo} again
\end{itemize}

\begin{tcolorbox}[breakable,colback=teal!4,colframe=teal!35!black,title={Before you move on}]
What does \mw{लेणो} mean? (\textquotedblleft{}To take\textquotedblright{}.) Say it once more.
\end{tcolorbox}
\section[\texorpdfstring{\mw{राखणो} (rākhṇo)}{राखणो (rākhṇo)}]{\mw{राखणो} (rākhṇo) — to keep, to put}

\label{lesson:MW-C53-rakhno}



\begin{quote}
Before the new one: say the Marwadi for to give, then the Marwadi for to take.
\end{quote}

\subsection*{You'll want to know: \mw{राखणो}}
\textbf{\mw{राखणो}} — \emph{rākhṇo} — \textquotedblleft{}to keep, to put\textquotedblright{}.

\begin{cousinweb}
\textbf{\mw{राखणो}} (\emph{rākhṇo}) is a tadbhava, a word inherited from Sanskrit and reshaped by sound change. It comes from \emph{rakṣati}, \textquotedblleft{}guards\textquotedblright{}. CDIAL lists Old Marwari \emph{rākhaï} under it (CDIAL 10547).
\end{cousinweb}

\begin{grammarlens}[title={What you've built}]
One more word for this chapter.
\end{grammarlens}

\subsection*{Your turn}
\begin{itemize}
\raggedright
  \item \emph{Say it:} \emph{rākhṇo}
  \item \emph{Say it:} \emph{rākhṇo}, once more
  \item \emph{Say it:} say \emph{rākhṇo}
  \item \emph{Recall:} say the Marwadi for to give, then the Marwadi for to take, then say \emph{rākhṇo} again
\end{itemize}

\begin{tcolorbox}[breakable,colback=teal!4,colframe=teal!35!black,title={Before you move on}]
What does \mw{राखणो} mean? (\textquotedblleft{}To keep, to put\textquotedblright{}.) Say it once more.
\end{tcolorbox}
